\chapter{Preludio y marco}

Tomo I fija el lenguaje fuerte del sistema.

\begin{observation}
En la arquitectura de publicación vigente, Tomo I funciona como volumen fundacional principal de MADELEINE.
\end{observation}

\begin{observation}
Su tarea es hacer visible, en una sola pieza editorial, el fundamento del sistema, su línea de lectura sobria y el núcleo estructural mínimo necesario para sostener el resto del corpus.
\end{observation}

\begin{observation}
El antiguo bloque denominado Base Axiomática queda absorbido, para fines de publicación, dentro de este Tomo I. Sus materiales pueden conservarse técnicamente como soporte interno, pero el volumen fundacional público del sistema es este tomo.
\end{observation}

\begin{observation}
Tomo II recibe desde aquí el canon intermedio de organización operativa y Tomo III recibe, desde Tomo II, la capa metodológica, curatorial y de proyección disciplinada.
\end{observation}

